% Frugal AI seminar, part 3: why India does not need to race for the frontier. The "panel view" items are
% opinions from an unnamed-speaker discussion and are labelled as such. Figures are from the Ref lines.

\section[Frontier]{Not Racing for the Frontier}

%%%%%%%%%%%%%%%%%%%%%%%%%%%%%%%%%%%%%%%%%%%%%%%%%%%
\begin{frame}[fragile]\frametitle{}
\begin{center}
{\Large Not Racing for the Frontier}

\medskip
{\small (A strategy, not a consolation prize)}
\end{center}
\end{frame}

%%%%%%%%%%%%%%%%%%%%%%%%%%%%%%%%%%%%%%%%%%%%%%%%%%%
\begin{frame}[fragile]\frametitle{The Frontier Race Is Expensive}
\begin{itemize}
\item The Frugal AI Hub's report argues frontier capability earns its cost ``in a narrow band'' of genuinely hard problems: advanced scientific research, long-horizon multi-step reasoning, and the cyber edge.
\item For the great majority of public-service and productivity problems it is ``not required, and is frequently the wrong tool''.
\item Owning the whole stack (compute, frontier models, data) is, it says, prohibitively expensive and, for nearly every country outside the US and China, unnecessary.
\item A panel view (opinion): being a second or third mover is acceptable, and even though the price per token is falling, the total bill is still large for many enterprises.
\end{itemize}
{\tiny (Ref: Frugal AI Hub, Minds, Models \& Markets, July 2026. The Hub's authors argue for frugal AI: read it as advocacy.)}
\end{frame}

%%%%%%%%%%%%%%%%%%%%%%%%%%%%%%%%%%%%%%%%%%%%%%%%%%%
\begin{frame}[fragile]\frametitle{A Labour-Rich Country}
\begin{itemize}
\item The Economic Survey 2023-24 estimates India needs about 7.85 million non-farm jobs a year until 2030. It is an official estimate, and others dispute its assumptions.
\item A panel view (opinion): not being at the frontier concentrates effort where India has strengths, in labour-intensive skills and services, and in AI that helps people rather than replaces them.
\item Small models fit that picture: cheap to deploy, close to the user, in the user's language.
\end{itemize}
{\tiny (Ref: Down To Earth, 22 July 2024, on the Economic Survey 2023-24)}
\end{frame}

%%%%%%%%%%%%%%%%%%%%%%%%%%%%%%%%%%%%%%%%%%%%%%%%%%%
\begin{frame}[fragile]\frametitle{Frontier-Minus-N: A Decision Rule}
\begin{itemize}
\item The Hub's report coins the term: a model ``deliberately built several generations behind the absolute frontier'', trading raw capability for lower cost, auditability, task fit and deployability under real connectivity and budget limits.
\item Its warning: it is only useful as a decision rule, not as an excuse after the fact. The questions are the task, its stakes, the available data, cost tolerance, auditability, connectivity, and how fast the capability changes.
\item The minister says AI's future will be shaped by ``smaller, efficient models'' and that nearly 95 per cent of AI work is done by small models. That is his statement, not a measurement; this talk does not rely on the figure.
\end{itemize}
{\tiny (Ref: Frugal AI Hub 2026; All India Radio, 21 January 2026)}
\end{frame}

%%%%%%%%%%%%%%%%%%%%%%%%%%%%%%%%%%%%%%%%%%%%%%%%%%%
\begin{frame}[fragile]\frametitle{Where Small Pays: Eye Screening}
\begin{itemize}
\item Google developed a model that detects diabetic retinopathy, a treatable cause of preventable blindness, from eye images. Aravind Eye Hospital was a foundational research partner.
\item By October 2024 it had supported more than 600,000 screenings; partners aim for 6 million AI-supported screenings over 10 years at no cost to patients, and were seeking local regulatory approvals.
\item It targets communities with ``a significant shortage of eye specialists''.
\item An imaging model, not a language model: the same frugal idea, a model fitted to one job, working at the point of care.
\end{itemize}
{\tiny (Ref: Google blog, 17 October 2024. The operator's own account.)}
\end{frame}
